\documentclass[../../../hsm_tok_q.tex]{subfiles}

\begin{document}
    \setcounter{figure}{0}
    Sさんのクラスでは，先生が示した問題をみんなで考えた．

    次の各問に答えよ．

    \begin{mycolorboxa}[title=［先生が示した問題］]

        \textsf{図\ref{fig/hsm_tok_2025_1st_02_01_q}}のように，円Oの円周を12等分する点に，%
        1から12までの自然数の番号を，小さい順で時計回りに付ける．

        1から12までの番号を付けた点のうち，2点を結んでできる線分が円Oの直径となるとき，%
        その2点を向かい合う点とする．

        例えば，1の点と7の点は向かい合う点である．

        \textsf{図\ref{fig/hsm_tok_2025_1st_02_01_q}}において，1組の向かい合う点を選び，%
        それぞれの点の番号のうち，小さい方の数を$a$，大きい方の数を$b$とする．

        $a$，$b$の平均値をA，$b^2-a^2$の値をBとするとき，BはAの何倍か求めなさい．
        \\
        \begin{minipage}{\linewidth}
            \vspace{.5\baselineskip}
            {\centering
                \setlength{\abovecaptionskip}{5pt}
                \includegraphics%
                    {\subfix{fig_hsm_tok_2025_1st_02/fig_hsm_tok_2025_1st_02_01_q.pdf}}
                    \captionof{figure}{} \label{fig/hsm_tok_2025_1st_02_01_q}
            }
        \end{minipage}
    \end{mycolorboxa}
    
    \begin{enumerate}
        \item ［先生が示した問題］で，BはAの\:\myboxf{　}\:倍と表すとき，%
            \:\myboxf{　}\:に当てはまる数を，次の{\textsf ア}〜\textsf{エ}のうちから%
            選び，記号で答えよ．
    
            \begin{tabularx}{\dimexpr\linewidth-\zw}{@{}LLLL@{}}
                {\textsf ア}\quad$3$ & {\textsf イ}\quad$4$ & {\textsf ウ}\quad$6$ & {\textsf エ}\quad$12$
            \end{tabularx}
    \end{enumerate}
    \newpage    
    Sさんのグループは，［先生が示した問題］をもとにして，次の問題を作った．

    \begin{mycolorboxa}[title=［Sさんのグループが作った問題］]

        \textsf{図\ref{fig/hsm_tok_2025_1st_02_02_q}}のように，円Oの円周を24等分する点に，%
        1から24までの自然数の番号を，小さい順で時計回りに付ける．

        1から24までの番号を付けた点のうち，2点を結んでできる線分が円Oの直径となるとき，%
        その2点を向かい合う点とする．

        \textsf{図\ref{fig/hsm_tok_2025_1st_02_02_q}}において，%
        異なる2組の向かい合う点を選び，1組目のそれぞれの点の番号のうち，%
        小さい方の数を$a$，大きい方の数を$b$とし，%
        2組目のそれぞれの点の番号のうち，小さい方の数を$c$，大きい方の数を$d$とする．

        $a$，$b$，$c$，$d$の平均値をP，$bd-ac$の値をQとするとき，%
        $\mathrm{Q}=24\,\mathrm{P}$となることを確かめてみよう．
        \\
        \begin{minipage}{\linewidth}
            \vspace{.5\baselineskip}
            {\centering
                \setlength{\abovecaptionskip}{5pt}
                \includegraphics%
                    {\subfix{fig_hsm_tok_2025_1st_02/fig_hsm_tok_2025_1st_02_02_q.pdf}}
                    \captionof{figure}{} \label{fig/hsm_tok_2025_1st_02_02_q}
            }
        \end{minipage}        
    \end{mycolorboxa}
    
    \begin{enumerate}[resume]
        \item ［Sさんのグループが作った問題］で，$\mathrm{Q}=24\,\mathrm{P}$となることを証明せよ．
    \end{enumerate}
\end{document}